% Effect of taking absolute values on the quadratic form on P_3.
\begin{tikzpicture}[x=1mm,y=1mm,font=\figurefont\fontsize{8}{10}\selectfont,>=Stealth,
  vertex/.style={circle,draw=FigureBlue,fill=FigureBlueFill,inner sep=0pt,
    minimum size=5mm,line width=.55pt},
  elabel/.style={font=\figurefont\fontsize{7.5}{9}\selectfont,FigureInk},
  vlabel/.style={font=\figurefont\fontsize{7.5}{9}\selectfont,FigureInk},
  panel/.style={font=\figurefont\fontsize{7.5}{9}\selectfont,text=FigureInk}]
  % Left panel: before absolute value
  \begin{scope}[xshift=0mm]
    \node[vertex] (l1) at (0,0) {$1$};
    \node[vertex] (l2) at (15,0) {$2$};
    \node[vertex] (l3) at (30,0) {$3$};
    \draw[FigureStroke,line width=.6pt] (l1) -- node[elabel,below=1mm] {\color{FigureRed}$-1$} (l2);
    \draw[FigureStroke,line width=.6pt] (l2) -- node[elabel,below=1mm] {\color{FigureRed}$-1$} (l3);
    \node[vlabel,above=1.5mm] at (l1) {$+1$};
    \node[vlabel,above=1.5mm] at (l2) {$-1$};
    \node[vlabel,above=1.5mm] at (l3) {$+1$};
    \node[vlabel] at (15,-8.5) {$\bm x^\top\bm A\bm x=-4$};
    \node[panel] at (15,-14.5) {(a) 赋值 $(1,-1,1)^\top$};
  \end{scope}
  % Arrow
  \draw[->,FigureStroke,line width=.65pt] (36,0) -- (46,0)
    node[midway,above,font=\figurefont\fontsize{7.5}{9}\selectfont] {$|\cdot|$};
  % Right panel: after absolute value
  \begin{scope}[xshift=52mm]
    \node[vertex] (r1) at (0,0) {$1$};
    \node[vertex] (r2) at (15,0) {$2$};
    \node[vertex] (r3) at (30,0) {$3$};
    \draw[FigureStroke,line width=.6pt] (r1) -- node[elabel,below=1mm] {\color{FigureBlue}$+1$} (r2);
    \draw[FigureStroke,line width=.6pt] (r2) -- node[elabel,below=1mm] {\color{FigureBlue}$+1$} (r3);
    \node[vlabel,above=1.5mm] at (r1) {$+1$};
    \node[vlabel,above=1.5mm] at (r2) {$+1$};
    \node[vlabel,above=1.5mm] at (r3) {$+1$};
    \node[vlabel] at (15,-8.5) {$|\bm x|^\top\bm A|\bm x|=+4$};
    \node[panel] at (15,-14.5) {(b) 取绝对值 $(1,1,1)^\top$};
  \end{scope}
\end{tikzpicture}
